




\subsection{ Toom's Contours}
\label{subsection:toricencoded}
We are about to zoom out from a single register encoded in the Toric code to describe a fully fault-tolerant computation. For a given quantum circuit, assumed to consist of gates $\mathbf{X}, \mathbf{Z}, \mathbf{CX}, \mathbf{H}, \mathbf{S}, \mathbf{T}$ and measurements, we define its Toric encoded circuit as its version where each qubit is encoded by the Toric code. Formally:
    \begin{definition}[Troic Encoded Circuit]
      Let $C$ be a quantum circuit. We say that $C$ is a \textit{Toric encoded circuit} if the following holds:
      \begin{enumerate}
        \item Any layer of the circuit consists of classical and quantum registers, such that there is a Toric code $\mathcal{C}$ such that the quantum registers are encoded by the \textit{nice encoding of} $\mathcal{C}$. At time $t=0$, the quantum registers are initialized to be either $\ket{\bar{0}}$, encoded $\ket{\mathbf{H}}$, encoded $\ket{\mathbf{S}}$, or encoded $\ket{\mathbf{T}}$. The classical registers are initialized to be zeros.
        \item The layers of $C$ on the quantum registers at even times are sweep-cycles, and the layers at odd times are realizations of either logical Paulis, $\mathbf{CX}$, or measurements in the computational and parity bases. On the classical registers, including the measured register, we allow any two-bit gates.
        \end{enumerate} 
    \end{definition}

For Toric encoded circuits, we define their \textit{base graph}, which can be thought of as a recording of a running from the Toric's point of view. The \textit{base graph} is obtained by taking the Toric's vertices at different times as the vertices and connecting any two sequential in time if they are support faces — qubits — which are either the same or are entangled by a quantum gate.
  \begin{definition}[Base graph $G_{o}$]
    Let $C$ be a Toric encoded circuit at depth $\mu$. Notice that measurements take quantum registers into classical, Thus any quantum register has a maximal depth, which we denote by $\{\mu_{i}$\}. Let's denote by $\{\mathcal{G}_{i}\}$ the skeleton of the $\infty$-extension of the Toric complexes that induce the codes. We define the base graph of $C$, denoted by $G_{o} = (V_{o}, E_{o})$, as follows:
    \begin{enumerate}
      \item The vertices $V_{o}$ are tuples, where the first coordinates correspond to a vertex in the disjoin union over the vertices in $\{\mathcal{G}_{i}\}$, and the second coordinate is associated with a time. Namely:
        \begin{equation*}
          \begin{split}
            V_{o} \colonequals \bigcup_{i} \bigl( V(\mathcal{G}_{i}) \times [\mu_{i}] \bigr)
          \end{split}
        \end{equation*}
      \item The edges $E_{o}$ are derived from the original edges in $\mathcal{G}_{1}, \ldots, \mathcal{G}_{k}$, with the addition that any two vertices in preceding times $t$ and $t+1$ that support qubits (faces) such that $C^{\infty}$ acts non-trivially on those qubits at time $t$ are also connected by an edge. Namely:
        \begin{equation*}
          \begin{split}
            E_{o} & \colonequals \big\{ \{(v,t), (u,t^{\prime}) \} :  \text{ if }     \\ 
              & \text{either there exists } i \text{ such }  \{v,u\} \in \mathcal{G}_{i} \text{ and } t^{\prime} = t+1 \\
            & \text{or there exists } u_i = (g_{i},l_{i}) \in C^{\infty} \text{ such }  v,u \in l_{i,2} \text{ and } t = l_{i,1}, t^{\prime} = l_{i,1}+1  \big\} 
          \end{split}
        \end{equation*}
Where we define a vertex $v$ to belong to location $l$ if there is a qubit-face $f$ in $l$ such that $v$ belongs to its associated face.
    \end{enumerate} 
  \end{definition}


As a concept, if the base graph is a recording of a run from the Torices' point of view, then a \textit{Toom's contour} is a recording of the deviation's evolution in the eyes of the Torices' vertices. We use the notion of Pauli-path to describe a deviation:




\begin{definition}

Let $P$ be a Pauli. We denote by $P^{\infty}$ the view of $P$ in the $\infty$-extension. We denote by $x(P^{\infty})$ the $d^{\prime}$-dimensional chain in the infinite complex that supports the $X$-component of $P^{\infty}$, and by $z(P^{\infty})$ the $d-d^{\prime}$-dimensional chain, such that the inversion of the $\phi$-transpose sends it to the $d^{\prime}$-dimensional chain that supports the $Z$-component of $P^{\infty}$, namely $\phi^{-1}(z(P^{\infty}))$ is the chain that supports the $Z$-component. 

%We define the reduced form of $P$ to be the Pauli obtained by the following procedure: for any finite connected chain that corresponds to a subset of the $X$ component of $P$'s support, substitute it with its reduced form (\Cref{def:reduced}). Then do the same for the $Z$ component in the dual picture. Observe that Paulis of a reduced Pauli-path act on the qubits in the $\infty$-extension.^i


\inb{\textbf{Consider removing.}} We extend the notion to Paulis in finite spaces (for example, Paulis on the circuit's register) by saying that a Pauli $P \in \mathcal{P}^{(\cdot)}$ is reduced if its view in the $\infty$-extension is reduced.
\end{definition}

\begin{definition}
  \label{def:paulipath}
  Consider a quantum circuit subjected to noise $C[\mathcal{E}]$, with width $w$ and depth $\mu$, whose registers are all encoded in the $(d,d^{\prime})$-Toric code with side length $n$ such that $C[\mathcal{E}]$ has an $\infty$-extension. We say that the set of Pauli $\{ (P_{i}) \}_{1}^{\mu+1}$, where for each $i$, $P_{i} \in \mathcal{P}^{[w]}$, is a Pauli-path for $C[\mathcal{E}]$. If $P_{1}$ is the identity and for any $t \in [\mu]$, the Pauli $P_{t+1}$ is defined such it satisfies the follow:  
  \begin{equation*}
    \begin{split}
      & P_{t+1} \in C|_{t} \circ P_{t} 
    \end{split}
  \end{equation*}
\inb{\textbf{Consider removing.}}  If there is a Pauli-path $P_{1},..,P_{\mu+1}$ and a time $t \in [\mu+1]$ such $P_{t}$ is infinite, then we say that time $t$ admits an infinite error. 
\end{definition}


For a Pauli-path $P_{1},..,P_{\mu+1}$, the Toom's contour is the restriction of the base graph to vertices that witness a deviation, namely, those that belong to qubits supporting $P_{1},P_{2},..,P_{\mu+1}$. 


\begin{definition}[Toom's Contour Graph of $C(\mathcal{E})$, (Standard form)]
    \label{def:contour}
    For a Toric encoded quantum circuit $C$, set of faults $\mathcal{E}$ and a Pauli-path $P_{1},..$, let us denote the base graph of $C[\mathcal{E}]$ by $G_{o}$. We define the \textit{Toom's contour graph} $\mathcal{H} = ( V_{\mathcal{H}}, E_{\mathcal{H}})$ as follows: 
    \begin{enumerate}
      \item The vertices $V_{\mathcal{H}} = V_{\mathcal{H}, X} \cup V_{\mathcal{H}, Z}$ correspond to vertices of faces supporting the Pauli-path. They consist of the triple $(v,t,\cdot)$ where $(v,t)$ is a vertex of the base graph, and the third is a symbolic coordinate that is either $\{X\}$ or $\{Z\}$ corresponding to the component of the Pauli that $v$ supports at time $t$.
        \begin{equation*}
          \begin{split}
            V_{\mathcal{H}, X} \colonequals & \big\{ v_{t} \times \{X\} : v_{t}=(v,t) \in V_{o} \\
            & \, \, \, \text{ and there exists a face } f \in \partial^{-} x(P_{t}) \text{ such } v \in f \big\} \\
             V_{\mathcal{H}, Z} \colonequals & \big\{ v_{t} \times \{Z\} : v_{t}=(v,t) \in V_{o} \\
             & \, \, \, \text{ and there exists a face } f \in \partial^{-} z(P_{t}) \text{ such } v \in f \big\} 
          \end{split}
        \end{equation*}

%We extend the action of potential walls to act on the vertices of $V_{\mathcal{H}}$ by acting on the associated points in $\mathbb{R}^{d}$. More concretely, for a vertex $v \in V_{\mathcal{H}}$, whose corresponding vertex on the infinite $\mathbb{R}^{d}$ Toric is $\tilde{v} \in \mathcal{G}$, and a potential wall $L_{i}$, the value of $L_{i}$ at $v$ is:
%        \begin{equation*}
%          \begin{split}
%            L_{i}(v) \colonequals L_{i}(\tilde{v})
%          \end{split}
%        \end{equation*}

      \item The edges $E_{\mathcal{H}}$ is partitioned into two types $A_{\mathcal{H}}$ and $F_{\mathcal{H}}$, where for any logical gate $g$ which applied in time $t$ in $C[\mathcal{E}]^{\infty} $ we add edge to $A_{\mathcal{H}}$ as follow: 
        \begin{enumerate}
          \item If $g$ is either a logical Pauli, a measurement in either the computational or the phase basis, a sweep-cycle, or an identity, then for any face $f$ in the register on which $g$ acts, we add:
            \begin{equation*}
              \begin{split}
                & \bigl\{  \{ \tilde{u} = (u,t, \cdot),  \tilde{v} = (v,t+1, \cdot)\} : u,v \in f \text{ and } \tilde{v},\tilde{u} \in V_{\mathcal{H},X} \\ 
                &  \, \, \, \text { or }  u,v \in \phi(f) \text{ and } \tilde{v},\tilde{u} \in V_{\mathcal{H},Z} \bigr\} 
              \end{split}
            \end{equation*}
            \item If $g$ is a logical $\mathbf{CX}$ then for any physical $\mathbf{CX}$ in its realization, between faces $f_{1}$ into $f_{2}$ at time $t$ we add: 
            \begin{equation*}
              \begin{split}
                & \bigl\{  \{ \tilde{u}  = (u,t, \cdot), \tilde{v} = (v,t+1, \cdot)\} :  u \in f_{1}\cup f_{2}, v \in f_{1}\cup f_{2}  \text{ and } \tilde{v},\tilde{u} \in V_{\mathcal{H},X}  \\
                &  \, \, \, \text { or } u \in \phi(f_{1})\cup \phi(f_{2}), v \in \phi(f_{1})\cup \phi(f_{2}) \text{ and } \tilde{v},\tilde{u} \in V_{\mathcal{H},Z} \bigr\} 
              \end{split}
            \end{equation*}
          \item If $g$ is a fault Pauli, then for any face $f$ in the register, we add:
            \begin{equation*}
              \begin{split}
                & \bigl\{  \{ \tilde{u} = (u,t, \cdot),  \tilde{v} = (v,t+1, \cdot)\} : u,v \in f \text{ , } \tilde{v},\tilde{u} \in V_{\mathcal{H},X} \text{, and } f \notin X(g)  \\ 
                &  \, \, \, \text { or }  u,v \in \phi(f) \text{ , } \tilde{v},\tilde{u} \in V_{\mathcal{H},Z} \text{, and } f \notin Z(g)  \bigr\} 
              \end{split}
            \end{equation*}
        \end{enumerate}
%        In addition for any fault $g$ of $\mathcal{E}^{\infty}$, that is applied at time $t$, and any face $f \in \text{Loc}(g)$, we add the edges:
%            \begin{equation*}
%              \begin{split}
%                 \bigl\{  \{(u,t),(v,t+1)\} : u,v \in f \text{ and } u,v \in V_{\mathcal{H},X} \bigr\} 
%              \end{split}
%            \end{equation*}
            To construct $F_{\mathcal{H}}$, we connect any pair of vertices $(v,t)$, $(u,t)$ if there exists a vertex $(w,t+1)$ such both of them are connected to it trough $A_{\mathcal{H}}$. Namely: 
        \begin{equation*}
          \begin{split}
            F_{\mathcal{H}} &= \big\{ e = \{u,v\} : \text{ exists } w \in V_{\mathcal{H}} \text{ such } \{v,w\},\{u,w\} \in A_{\mathcal{H}} \big\} 
          \end{split}
        \end{equation*}
    \end{enumerate}
    We call $A_{\mathcal{H}}$ and $F_{\mathcal{H}}$ the \textit{arrows} and \textit{forks} edges respectively. 
\end{definition}


%\inb{\textbf{Insert an example.}}




\begin{remark}
One might expect to have a differentiation between measurement in the computational and the phases basis. Namely, a rule of connecting arrows of the form:
\begin{center}
\begin{minipage}{0.75\textwidth}
  \textit{
            If $g$ is a measurement in the phases base, then for any face $f$ in the register on which $g$ acts, we add:
            \begin{equation*}
              \begin{split}
                & \bigl\{  \{ \tilde{u} = (u,t, \cdot),  \tilde{v} = (v,t+1, \cdot)\} : u \in f , v \in \phi(f) \text{ and } \tilde{u} \in V_{\mathcal{H},X}, \tilde{v} \in V_{\mathcal{H},Z}, \\ 
                &  \, \, \, \text { or }  u\in \phi(f), v \in f \text{ and } \tilde{u} \in V_{\mathcal{H},Z} , \tilde{v} \in  V_{\mathcal{H},X} \bigr\} 
              \end{split}
            \end{equation*}
  }
\end{minipage}
\end{center}
Such intuition is right in the sense that this is the right way to keep track of the faults. Yet, later we will state \Cref{claim:replace} whose corollary is that instead of analyzing the running of the Toric encoded circuit, it is sufficient to analyze the running of a 'simpler' circuit which doesn't contain measurement in the phase base. In particular, the claim considers the circuit obtained from a given Toric encoded circuit by replacing each of the classical decoding procedures with noisy sweep-rule cycles. Then it states that to derive an upper bound on the failure probability of the original circuit, it's enough to bound the probability that the modified circuit does not admit an infinite 'time-connected cluster' (slice, \Cref{def:slice}).

\end{remark}




\begin{definition}
  \label{def:paulipathupstep}
With the same notations as in \Cref{def:paulipath}, let $\tau$ be a step in $C[\mathcal{E}]$. Recall that each gate in $C[\mathcal{E}]$ acts on at most $\Delta$ qubits. Let $\Delta^{\prime} \in [\Delta]$. In addition, denote by $g$ and $t$ the gate and its time which is applied at step $\tau$. We say that the set of Pauli $P_{1},\ldots,P_{t},P_{\star}$ is a Pauli-path up to step $(\tau,\Delta^{\prime})$, if the $P_{1},\ldots,P_{t}$ is a Pauli-path for the circuit obtained from $C[\mathcal{E}]$ by taking only its first $t-1$ layers. And $P_{\star}$ satisfies the following: let $g^{\prime}_{1},\ldots,g_{l}^{\prime} = g$ be the gates applied at time $t$ at step up to, and including, $\tau$, placed by order. Let $P^{\prime}_{0},P^{\prime}_{\Delta} \in \mathcal{P}^{(\cdot)}$ be Paulis such:
  \begin{equation*}
    \begin{split}
      P_{0}^{\prime} \in \bigl(g^{\prime}_{l-1}g^{\prime}_{l-2}..g_{1}\bigr) \circ P_{t} \quad P^{\prime}_{\Delta} \in g \circ P_{0}^{\prime}
    \end{split}
  \end{equation*}
Now, since $P^{\prime}_{\Delta}$ is at a Hamming distance of at most $\Delta$ from $P_{0}^{\prime}$, there is a set of Paulis:
  \begin{equation*}
    \begin{split}
      P_{0}^{\prime},P_{1}^{\prime},P_{2}^{\prime},..,P^{\prime}_{\Delta^{\prime}},..,P^{\prime}_{\Delta}
    \end{split}
  \end{equation*}
  Then set $P_{\star} = P^{\prime}_{\Delta^{\prime}}$. 
\end{definition}


\inb{\textbf{Rewrite. }} For the analyses, we would like to extend Toom's contour to record in their final layer a partial time-step, namely a computation up to a certain step \Cref{def:steps}. For this, we substitute in the definition above the Pauli-path with a Pauli-path up to a time \inb{\textbf{Prefix?}}. Moreover, since we would like to refer to the connected chain at the final layer as the 'final error clusters', we add at the end a layer whose purpose is to connect in time (by arrows) the connected components of Toom's contour in the prior layer, which are connected by checks (their union is identified with a connected chain).
\begin{definition}[Analysis Toom's Contour Up to Step $(\tau,\Delta^\prime)$]
For a Toric encoded quantum circuit $C$, a set of faults $\mathcal{E}$, a step $\tau$ according to the standard order, and an integer $\Delta^{\prime} \in [\Delta]$, we define the \textit{analysis Toom's contour up to step $(\tau,\Delta^\prime)$}, denoted by $\mathcal{H}_{\tau,\Delta^{\prime}} = (V_{\mathcal{H}}, E_{\mathcal{H}})$, to be the graph obtained by taking a Pauli-path up to step $(\tau,\Delta^{\prime})$ $P_{1}, \ldots, P_{t}, P_{\star}$, and gluing to the Toom's contour induced by the prefix $P_{1}, \ldots, P_{t}$ a layer defined by adding vertices and edges as follows:
  \begin{enumerate}
    \item We add to the vertices in $V_{\mathcal{H},\{X\}}$ the vertices that support $P_{\star}$, namely the subset:
      \begin{equation*}
          \begin{split}
            \bigl\{ v_{t+1}  & : v_{t+1}=( v,t+1,\{X\}  ) \text{ if there exists a face } f \in \partial^{-}x(P_{\star}) \\ 
            & \text{ such } v \in f \in \partial^{-}x(P_{\star}) \bigr\} 
          \end{split}
        \end{equation*}
        Similarly we add to $V_{\mathcal{H}, Z}$ the subset: 
      \begin{equation*}
          \begin{split}
            \bigl\{ v_{t+1}  & : v_{t+1}=(  v,t+1,\{Z\} ) \text{ if there exists a face } f \in \partial^{-} z(P_{\star}) \\
            & \text{ such } v \in f \bigr\} 
          \end{split}
        \end{equation*}
      \item For the edges, we add the following arrows:
        \begin{equation*}
          \begin{split}
            \bigl\{ \bigl\{ \tilde{u} & = (u,t+1 -1, \cdot ),  \tilde{v}= (v,t+1,\cdot)  \bigr\}  \text{ if  either there exists }\\ 
           & \text{ a } d^{\prime}-2  \text{ dimensional face } f  \text{such } \tilde{u},\tilde{v} \in V_{\mathcal{H},X} \text{ and } u,v \in f  \\
          & \text{or there exists a } d - d^{\prime}-2  \text{ dimensional face } f \text{ such } \tilde{u},\tilde{v} \in V_{\mathcal{H},Z} \text{ and } v,u \in f   \bigr\}  
          \end{split}
        \end{equation*}
      \item Finally, we add forks for vertices at time $t+1$ to preserve the property that any two vertices with the same time coordinate, which are connected to a third (with a $+1$ higher time coordinate) via arrows, are also connected by a fork.
  \end{enumerate}


\end{definition}




